\documentclass[11pt]{article}
\usepackage{worksheet_style}

% ---- Toggle: uncomment for filled version ----
%\worksheetfilledtrue
\begin{document}
\worksheetheader{8}{Partial Derivatives}

\begin{background}
Now we differentiate. The idea is simple: hold all variables constant except one, and differentiate with respect to that one. This gives us partial derivatives---the multivariate analog of $f'(x)$. Today we define them, compute them, and see how they measure the slope of a surface in each coordinate direction.
\end{background}

%=============================================================
\wsection{Partial Derivatives}

\begin{defbox}{Partial Derivative (Limit Definition)}
The \textbf{partial derivative} $f_x(a,b)$ is the instantaneous rate of change of $f(x,y)$ in the $x$-direction at the point $(a,b)$, with $y$ held fixed at $b$. Geometrically, slice the graph of $f$ with the vertical plane $y=b$; the resulting cross-section is a single-variable curve $z=f(x,b)$, and $f_x(a,b)$ is the \textbf{slope of the tangent line} to that curve at $x=a$. In other words, partial derivatives are ordinary single-variable slopes taken along axis-aligned slices of the surface.
\[f_x(x,y) = \frac{\partial f}{\partial x} = \wbml{\lim_{h\to 0} \frac{f(x+h, y) - f(x,y)}{h}}\]

With respect to $y$:
\[f_y(x,y) = \frac{\partial f}{\partial y} = \wbml{\lim_{h\to 0} \frac{f(x, y+h) - f(x,y)}{h}}\]

\textbf{Practical rule:} \wb{To find $f_x$, treat $y$ as a constant and differentiate with respect to $x$. For $f_y$, treat $x$ as constant.}
\end{defbox}

\wsubsection{Geometric Interpretation}

-- $f_x(a,b)$ is the slope of the tangent line to $z = f(x,b)$ at $x = a$

-- Measures rate of change of $f$ in the $x$-direction while $y$ is held fixed

-- Similarly, $f_y(a,b)$ is the slope of $z = f(a,y)$ at $y = b$

\begin{center}
\begin{tikzpicture}[scale=0.7]
  \begin{axis}[
      view={-60}{30},
      axis lines=center,
      xlabel={$x$},
      ylabel={$y$},
      zlabel={$z$},
      width=12cm,
      height=10cm,
      ticks=none,
    ]

    % Plot the surface z = x^2 + y^2
    \addplot3[
      surf,
      opacity=0.4,
      samples=30,
      samples y=30,
      domain=-2:2,
      y domain=-2:2,
      colormap/viridis
    ]
    {x^2 + y^2};

    % Tangent line at point (1, 1, 2)
    % Equation: z = 2 + 2(x - 1) + 2(y - 1)
    \addplot3[
      thick,
      color=bcred,
      domain=0:2
    ]
    (1 + x, 1 + x, {2 + 2*x + 2*x});

    % Projection of tangent line onto the xz-plane (y = constant)
    \addplot3[
      thick,
      dashed,
      color=bcblue,
      domain=0:2
    ]
    (1 + x, 1, {2 + 2*x});

    % Projection of tangent line onto the yz-plane (x = constant)
    \addplot3[
      thick,
      dashed,
      color=bcgreen,
      domain=0:2
    ]
    (1, 1 + x, {2 + 2*x});

    % Mark the point of tangency
    \addplot3[
      only marks,
      mark=*,
      mark size=3pt,
      color=black
    ]
    coordinates {(1, 1, 2)};

  \end{axis}
\end{tikzpicture}

{\small Tangent line (red) at $(1,1,2)$ on $z = x^2+y^2$. Blue dashed: $xz$-projection (slope $= f_x = 2$). Green dashed: $yz$-projection (slope $= f_y = 2$).}
\end{center}

\begin{example}[Using the Limit Definition]
Use the limit definition to find $f_x$ for $f(x,y) = x^2 y + 3xy^2$.

\wbbox{
\textbf{1. Set up:}
$f_x = \lim_{h\to 0}\frac{(x+h)^2 y + 3(x+h)y^2 - x^2 y - 3xy^2}{h}$

\textbf{2. Expand:}
$= \lim_{h\to 0}\frac{x^2 y + 2xhy + h^2 y + 3xy^2 + 3hy^2 - x^2 y - 3xy^2}{h}$

$= \lim_{h\to 0}\frac{2xhy + h^2 y + 3hy^2}{h}$

\textbf{3. Simplify and evaluate:}
$= \lim_{h\to 0}(2xy + hy + 3y^2) = \boxed{2xy + 3y^2}$
}{4.5cm}
\end{example}

\begin{example}[Computing Partial Derivatives]
Find all first partial derivatives:
\begin{enumerate}[label=(\alph*)]
  \item $f(x,y) = x^3 e^y + \sin(xy)$

  \wbbox{
  $f_x = 3x^2 e^y + y\cos(xy)$

  $f_y = x^3 e^y + x\cos(xy)$
  }{2cm}

  \item $g(x,y) = \dfrac{x^2}{x+y}$

  \wbbox{
  $g_x = \frac{2x(x+y) - x^2}{(x+y)^2} = \frac{x^2+2xy}{(x+y)^2}$ \quad (quotient rule in $x$)

  $g_y = \frac{0 - x^2}{(x+y)^2} = \frac{-x^2}{(x+y)^2}$
  }{2.5cm}

  \item $h(x,y,z) = xye^z + z\ln(x+y)$

  \wbbox{
  $h_x = ye^z + \frac{z}{x+y}$, \quad $h_y = xe^z + \frac{z}{x+y}$, \quad $h_z = xye^z + \ln(x+y)$
  }{2cm}
\end{enumerate}
\end{example}

%=============================================================
\wsection{Derivative Rules for Partials}

\begin{formulabox}{Rules (Hold Other Variables Constant)}
\begin{itemize}[nosep, leftmargin=*]
  \item \textbf{Linearity:} $\partial_x[af + bg] = af_x + bg_x$
  \item \textbf{Product:} $\partial_x[fg] = \wbm{f_x g + f g_x}$
  \item \textbf{Quotient:} $\partial_x[f/g] = \wbm{(f_x g - fg_x)/g^2}$
  \item \textbf{Chain:} $\partial_x[f(g(x,y))] = \wbm{f'(g)\cdot g_x}$
  \item \textbf{Zero partial:} if $f_x = 0$ everywhere $\implies$ $f$ does not depend on $x$
  \item \textbf{Constants:} $\partial_x[c] = 0$; treat other variables as constants
\end{itemize}
\end{formulabox}

%=============================================================
\wsection{Higher-Order Partial Derivatives}

\begin{defbox}{Second-Order Partial Derivatives}
A \textbf{second-order partial derivative} is a partial derivative of a partial derivative. The two \emph{pure} second partials $f_{xx}$ and $f_{yy}$ measure \textbf{concavity along coordinate directions} — how fast the slope $f_x$ changes as you move further in $x$, and similarly for $y$ — exactly analogous to the second derivative of a single-variable function. The two \emph{mixed} partials $f_{xy}$ and $f_{yx}$ measure something new: how the slope in \emph{one} direction changes as you move in the \emph{other} direction, capturing the graph's twist or saddle-like behavior at a point.
\begin{align*}
  f_{xx} &= \frac{\partial^2 f}{\partial x^2} = \frac{\partial}{\partial x}\left(\frac{\partial f}{\partial x}\right) \quad \text{(pure, in $x$)}\\[4pt]
  f_{yy} &= \frac{\partial^2 f}{\partial y^2} \quad \text{(pure, in $y$)}\\[4pt]
  f_{xy} &= \frac{\partial^2 f}{\partial y\,\partial x} = \frac{\partial}{\partial y}\left(\frac{\partial f}{\partial x}\right) \quad \text{(mixed: $x$ first, then $y$)}\\[4pt]
  f_{yx} &= \frac{\partial^2 f}{\partial x\,\partial y} \quad \text{(mixed: $y$ first, then $x$)}
\end{align*}
\end{defbox}

\begin{thmbox}{Clairaut's Theorem (Equality of Mixed Partials)}
If $f_{xy}$ and $f_{yx}$ are both continuous on an open region, then:
\[\wbml{f_{xy} = f_{yx}}\]
The order of differentiation does not matter (under mild continuity conditions).
\end{thmbox}

\wsubsection{Why Clairaut's Theorem Works (Proof Outline)}

-- Consider the second-order difference quotient over a small rectangle $[x, x+h]\times[y, y+k]$

-- Define $\Delta = f(x+h,y+k) - f(x+h,y) - f(x,y+k) + f(x,y)$

-- Apply the Mean Value Theorem twice: first in $x$, then in $y$ $\implies$ $\Delta = f_{xy}(c_1, d_1)\cdot hk$

-- Apply MVT in the other order: first in $y$, then in $x$ $\implies$ $\Delta = f_{yx}(c_2, d_2)\cdot hk$

-- As $h, k \to 0$, continuity gives $f_{xy}(x,y) = f_{yx}(x,y)$

\wsubsection{When Clairaut Fails -- A Counterexample}

-- Consider $f(x,y) = \begin{cases} \frac{xy(x^2-y^2)}{x^2+y^2} & (x,y)\neq(0,0) \\ 0 & (x,y)=(0,0)\end{cases}$

-- One can show $f_{xy}(0,0) = -1$ and $f_{yx}(0,0) = 1$ -- the mixed partials are \textit{not equal} at the origin because they are not continuous there

\begin{example}[Clairaut Verification]
Let $f(x,y) = x^3 y^2 + e^{xy}$. Find all four second-order partials and verify $f_{xy} = f_{yx}$.

\wbbox{
\textbf{1. First partials:}

$f_x = 3x^2 y^2 + ye^{xy}$, \quad $f_y = 2x^3 y + xe^{xy}$

\textbf{2. Second partials:}

$f_{xx} = 6xy^2 + y^2 e^{xy}$

$f_{yy} = 2x^3 + x^2 e^{xy}$

$f_{xy} = \frac{\partial}{\partial y}[3x^2 y^2 + ye^{xy}] = 6x^2 y + e^{xy} + xye^{xy}$

$f_{yx} = \frac{\partial}{\partial x}[2x^3 y + xe^{xy}] = 6x^2 y + e^{xy} + xye^{xy}$

\textbf{3. Result:} $f_{xy} = f_{yx}$ \;\checkmark
}{5.5cm}
\end{example}

\begin{example}[Application: Container Design]
A cylindrical container has volume $V(r,h) = \pi r^2 h$ and surface area $S(r,h) = 2\pi r h + 2\pi r^2$. If $r = 5$ cm and $h = 10$ cm, how does $V$ change if we increase $r$ by $0.1$ cm? If we increase $h$ by $0.1$ cm?

\wbbox{
\textbf{1. Compute partials of $V$:}

$V_r = 2\pi r h = 2\pi(5)(10) = 100\pi \approx 314$ cm$^3$/cm

$V_h = \pi r^2 = 25\pi \approx 78.5$ cm$^3$/cm

\textbf{2. Estimate changes:}

$\Delta r = 0.1$: $\Delta V \approx V_r \cdot 0.1 = 10\pi \approx 31.4$ cm$^3$

$\Delta h = 0.1$: $\Delta V \approx V_h \cdot 0.1 = 2.5\pi \approx 7.85$ cm$^3$

\textbf{3. Result:} Increasing radius has $\sim 4\times$ more effect on volume than increasing height.
}{5cm}
\end{example}

%=============================================================
\begin{practice}

\prob
Find $f_x$ and $f_y$ for $f(x,y) = x^2\cos y + y\ln x$.

\wans{$f_x = $ \wbml{$2x\cos y + y/x$}, \; $f_y = $ \wbml{$-x^2\sin y + \ln x$}}

\vspace{2cm}

\prob
Find all first partials of $f(x,y,z) = xyz + z^2 \sin(x)$.

\wbbox{
$f_x = yz + z^2\cos x$, \quad $f_y = xz$, \quad $f_z = xy + 2z\sin x$
}{2cm}

\prob
For $f(x,y) = \arctan(y/x)$, show that $f_{xx} + f_{yy} = 0$ (i.e., $f$ is \textbf{harmonic}).

\wbbox{
$f_x = \frac{-y/x^2}{1+y^2/x^2} = \frac{-y}{x^2+y^2}$, \quad $f_y = \frac{1/x}{1+y^2/x^2} = \frac{x}{x^2+y^2}$

$f_{xx} = \frac{2xy}{(x^2+y^2)^2}$, \quad $f_{yy} = \frac{-2xy}{(x^2+y^2)^2}$

$f_{xx}+f_{yy} = 0$ \;\checkmark
}{4cm}

\prob
Use the limit definition to compute $f_x(0,0)$ for
$f(x,y) = \begin{cases} \frac{xy}{x^2+y^2} & (x,y)\neq(0,0) \\ 0 & (x,y)=(0,0)\end{cases}$

\wbbox{
$f_x(0,0) = \lim_{h\to 0}\frac{f(h,0)-f(0,0)}{h} = \lim_{h\to 0}\frac{0-0}{h} = \boxed{0}$
}{2.5cm}

\prob
Find $f_{xxy}$ for $f(x,y) = e^x \sin(xy)$.

\wbbox{
$f_x = e^x\sin(xy) + ye^x\cos(xy)$

$f_{xx} = e^x\sin(xy) + ye^x\cos(xy) + ye^x\cos(xy) - y^2 e^x\sin(xy)$
$\quad = e^x[(1-y^2)\sin(xy) + 2y\cos(xy)]$

$f_{xxy} = e^x[-2y\sin(xy) + x(1-y^2)\cos(xy) + 2\cos(xy) - 2xy\sin(xy)]$
}{5cm}

\end{practice}

\end{document}
